\usepackage{mathrsfs}
\let\mathcal\mathscr
\multlinegap0pt

\newcommand{\citehuybrechts}{\cite{Huy99}}
\begin{DefTralics}
\newcommand{\citehuybrechts}{[Huy99]}
\end{DefTralics}

\newcommand{\Changel}{\mathcode`l="7160}
\newcommand{\Changelback}{\mathcode`l="716C}
\newcommand{\NoChangel}[1]{%
\expandafter\let\csname old\string#1\endcsname=#1
\let#1=\relax
\newcommand{#1}{\mathcode`l="716C\csname old\string#1\endcsname\mathcode`l="7160 }%
}

\NoChangel{\log}
\NoChangel{\ln}
\NoChangel{\lim}
\NoChangel{\limsup}
\NoChangel{\liminf}
\NoChangel{\varlimsup}
\NoChangel{\varliminf}
\NoChangel{\varinjlim}
\NoChangel{\varprojlim}


\newcommand{\Psfrac}[2]{(\sfrac{#1}{#2})}
\newcommand{\psfrac}[2]{\sfrac{(#1)}{#2}}
\newcommand{\spfrac}[2]{\sfrac{#1}{(#2)}}
\newcommand{\pspfrac}[2]{\sfrac{(#1)}{(#2)}}
\newcommand{\sbullet}{{\scriptscriptstyle\bullet}}
\newcommand\mto{\mathchoice{\longmapsto}{\mapsto}{\mapsto}{\mapsto}}
\newcommand\hto{\mathrel{\lhook\joinrel\to}}

\renewcommand\thesubsection{\thesection.\Alph{subsection}}

\theoremstyle{plain}
\newtheorem{theorem}{Theorem}[section]
\newtheorem{conjecture}[theorem]{Conjecture}
\newtheorem{corollary}[theorem]{Corollary}
\newtheorem{lemma}[theorem]{Lemma}
\newtheorem{proposition}[theorem]{Proposition}

\theoremstyle{definition}
\newtheorem{definition}[theorem]{Definition}
\newtheorem{remark}[theorem]{Remark}
\newtheorem{example}[theorem]{Example}

\usepackage[all,cmtip]{xy}\let\labelstyle\textstyle


\newcommand{\R}{\ensuremath{\mathbb{R}}}
\newcommand{\Q}{\ensuremath{\mathbb{Q}}}
\newcommand{\Z}{\ensuremath{\mathbb{Z}}}
\newcommand{\CC}{\ensuremath{\mathbb{C}}}\let\C\CC
\newcommand{\N}{\ensuremath{\mathbb{N}}}
\newcommand{\PP}{\ensuremath{\mathbb{P}}}

\newcommand{\holom}[3]{\ensuremath{#1:#2\to #3}}
\newcommand{\fibre}[2]{\ensuremath{#1^{-1} (#2)}}

\newcommand\sF{{\mathcal F}}
\newcommand\sO{{\mathcal O}}

\newcommand{\chow}[1]{\ensuremath{\mathcal{C}(#1)}}

\newcommand{\barS}{\ensuremath{\overline{S}}}
\newcommand{\blS}{\ensuremath{\widetilde{S \times S}}}

\DeclareMathOperator{\Def}{Def}
\DeclareMathOperator{\Eff}{Eff}
\DeclareMathOperator{\id}{id}
\DeclareMathOperator{\Nef}{Nef}
\DeclareMathOperator{\NS}{NS}
\DeclareMathOperator{\pic}{Pic}
\DeclareMathOperator{\supp}{Supp}
\DeclareMathOperator{\Td}{Td}
\DeclareMathOperator{\Todd}{Todd}
\DeclareMathOperator{\tr}{tr}

